% AUTO-GENERATED by scripts/exp_real_input_40_event_time_rescaling_error_budget.py
\begin{table}[H]
\centering
\scriptsize
\setlength{\tabcolsep}{6pt}
\renewcommand{\arraystretch}{1.15}
\caption{分位数调整后的事件时间相关尺度：$\tau_{\mathrm{eff}}^{(e)}(k)$ 与扭曲混合时间 $\tau_{\mathrm{mix}}$ 的对比（真实输入 40 状态核；$t=0$ Parry 平衡态）。设步数时间指数包络为 $\rho^n$，其中 $\rho=\rho_{\mathrm{corr}}(0)=0.618033988750$，$\tau_{\mathrm{corr}}(0)=1/(-\log\rho)=2.078086921235$。对事件阈值 $k$ 的击中时间 $\tau(k)$，以其精确分位数 $n_{\mathrm{lo}},n_{\mathrm{med}},n_{\mathrm{hi}}$ 诱导事件时间的有效相关尺度 $\tau_{\mathrm{eff}}^{(e)}(k):=(k/\tau(k))\,\tau_{\mathrm{corr}}(0)$。表中给出 $[q_{\mathrm{lo}},q_{\mathrm{hi}}]$ 区间内的尺度区间 $[\tau_{\mathrm{eff,lo}}^{(e)}(k),\tau_{\mathrm{eff,hi}}^{(e)}(k)]=[(k/n_{\mathrm{hi}})\tau_{\mathrm{corr}}(0),(k/n_{\mathrm{lo}})\tau_{\mathrm{corr}}(0)]$，其中 $q_{\mathrm{lo}}=0.05$，$q_{\mathrm{hi}}=0.95$，并比较 $\tau_{\mathrm{mix}}=19.747340579192$ 与该区间的比值。}
\label{tab:real-input-40-event-time-tau-eff-vs-tau-mix}
\begin{tabular}{r r r r r r}
\toprule
$k$ & $k/\mu_e$ & $\tau_{\mathrm{eff,lo}}^{(e)}(k)$ & $\tau_{\mathrm{eff,med}}^{(e)}(k)$ & $\tau_{\mathrm{eff,hi}}^{(e)}(k)$ & $\tau_{\mathrm{mix}}/\tau_{\mathrm{eff,hi}}^{(e)} \sim \tau_{\mathrm{mix}}/\tau_{\mathrm{eff,lo}}^{(e)}$\\
\midrule
6 & 16.332032046126 & 0.623426076371 & 0.831234768494 & 1.039043460618 & $19.005307600373\,\sim\,31.675512667289$\\
7 & 19.054037387147 & 0.632461236898 & 0.808144913814 & 1.039043460618 & $19.005307600373\,\sim\,31.223005343470$\\
8 & 21.776042728168 & 0.639411360380 & 0.791652160470 & 0.977923257052 & $20.193139325397\,\sim\,30.883624850607$\\
9 & 24.498048069189 & 0.644923527280 & 0.779282595463 & 0.984356962690 & $20.061158022616\,\sim\,30.619662245046$\\
10 & 27.220053410210 & 0.649402162886 & 0.799264200475 & 0.944584964198 & $20.905838360411\,\sim\,30.408492160597$\\
\bottomrule
\end{tabular}
\end{table}
